\chapter{Diseases of the Skin}
\label{ch:skin}
The skin is the largest organ of the body, covering approximately 1.7 m$^2$ and comprising about 16\% of body weight. It consists of three layers: the epidermis (keratinocytes, melanocytes, Langerhans cells), dermis (collagen, elastin, blood vessels, nerves, hair follicles, sweat glands), and hypodermis (subcutaneous fat). The skin serves as a barrier against pathogens, ultraviolet radiation, dehydration, and mechanical injury, and plays roles in thermoregulation, sensation, and immune surveillance. Dermatological diseases are among the most common reasons for medical consultation.

%-------------------------------------------------------------
\section{Acne and Related Disorders}
%-------------------------------------------------------------

%-------------------------------------------------------------

%-------------------------------------------------------------

%-------------------------------------------------------------

%-------------------------------------------------------------

%-------------------------------------------------------------

\label{sec:Acne and Related Disorders}
%-------------------------------------------------------------%-------------------------------------------------------------

\subsection{Acne Vulgaris}
\label{sec:Acne Vulgaris}
Acne vulgaris is a chronic inflammatory disorder of the pilosebaceous unit, affecting up to 85\% of adolescents and young adults. The condition results from four pathogenic factors: follicular hyperkeratinization, increased sebum production (androgen-mediated), \textit{Cutibacterium acnes} (formerly \textit{Propionibacterium acnes}) colonization, and inflammation. Clinical lesions include open comedones (blackheads), closed comedones (whiteheads), papules, pustules, nodules, and cysts; acne predominantly affects the face, chest, and back. Severity is classified as mild (comedonal), moderate (papulopustular), or severe (nodulocystic, conglobata). Diagnosis is clinical. Management depends on severity: topical retinoids (tretinoin, adapalene) and benzoyl peroxide for mild acne; combined topical therapy with topical antibiotics (clindamycin) for moderate acne; oral antibiotics (doxycycline, minocycline), hormonal therapy (combined oral contraceptives, spironolactone for females), and isotretinoin (vitamin A derivative, the only therapy that addresses all four pathogenic factors) for severe acne. Isotretinoin requires iPLEDGE program enrollment due to teratogenicity. Prognosis is generally favorable; many patients improve with age.

\subsection{Rosacea}
\label{sec:Rosacea}
Rosacea is a chronic inflammatory facial dermatosis primarily affecting middle-aged adults with fair skin. The condition results from innate immune dysregulation, neurovascular dysfunction, and \textit{Demodex folliculorum} mite overgrowth; triggers include sunlight, alcohol, spicy foods, hot beverages, and emotional stress. Clinical features present with four subtypes: erythematotelangiectatic (persistent facial erythema, telangiectasia, flushing), papulopustular (erythema with papules and pustules, resembling acne), phymatous (skin thickening, particularly rhinophyma), and ocular (blepharitis, conjunctivitis, keratitis). Diagnosis is clinical. Management includes sun protection, topical metronidazole or azelaic acid, topical ivermectin, oral doxycycline (low-dose, anti-inflammatory), and laser therapy for telangiectasia; severe rhinophyma may require surgical treatment. Prognosis is chronic with flares; treatment provides control but not cure.

%-------------------------------------------------------------
\section{Autoimmune Blistering Diseases}
%-------------------------------------------------------------

%-------------------------------------------------------------

%-------------------------------------------------------------

%-------------------------------------------------------------

%-------------------------------------------------------------

%-------------------------------------------------------------

\label{sec:Autoimmune Blistering Diseases}
%-------------------------------------------------------------%-------------------------------------------------------------

\subsection{Bullous Pemphigoid}
\label{sec:Bullous Pemphigoid}
Bullous pemphigoid (BP) is the most common autoimmune blistering disease, predominantly affecting the elderly (peak incidence 70--80 years). The condition results from IgG autoantibodies targeting BP180 and BP230, components of hemidesmosomes at the dermal-epidermal junction, causing subepidermal blister formation. Clinical features include tense blisters on an erythematous or urticarial base, pruritus, and predominance on the trunk and flexural surfaces; mucosal involvement is uncommon (10--30\%). Diagnosis involves skin biopsy showing subepidermal split with eosinophilic infiltrate, direct immunofluorescence showing linear IgG and C3 at the basement membrane zone, and anti-BP180/BP230 antibodies by ELISA. Management includes potent topical corticosteroids (clobetasol, first-line for localized and generalized disease), systemic corticosteroids for extensive disease, and steroid-sparing agents (doxycycline, azathioprine, mycophenolate, rituximab). Prognosis is generally favorable but may be chronic.

\subsection{Dermatitis Herpetiformis}
\label{sec:Dermatitis Herpetiformis}
Dermatitis herpetiformis (DH) is the cutaneous manifestation of celiac disease, characterized by intensely pruritic, grouped vesicles and papules on an erythematous base, symmetrically distributed on extensor surfaces (elbows, knees, buttocks, back). The condition results from IgA antibody deposition at the dermal papillae, triggered by gluten ingestion; nearly all patients have underlying celiac disease, though gastrointestinal symptoms may be absent. Clinical features include intensely pruritic, grouped vesicles on extensor surfaces. Diagnosis is based on skin biopsy showing granular IgA deposition at the dermal papillae. Management includes a strict gluten-free diet and dapsone (which rapidly controls skin symptoms); sulfapyridine is an alternative for patients intolerant to dapsone. Prognosis is excellent with dietary compliance.

\subsection{Epidermolysis Bullosa}
\label{sec:Epidermolysis Bullosa}
Epidermolysis bullosa (EB) is a group of inherited blistering disorders characterized by mechanical fragility of the skin and mucous membranes, with an overall prevalence of 1 in 50,000. The condition results from mutations in genes encoding structural proteins of the skin, classified by the level of blister formation within the basement membrane zone: EB simplex (EBS, keratin 5/14 mutations, intraepidermal cleavage), junctional EB (JEB, laminin 332/collagen XVII mutations, lamina lucida cleavage), and dystrophic EB (DEB, collagen VII mutations, sublamina densa cleavage). Clinical features include blister formation at sites of minor friction or trauma, ranging from mild acral blistering (EBS) to widespread blistering, scarring, milia, nail dystrophy, and mitten-hand deformities from repeated scarring (severe DEB). Extracutaneous involvement includes oral, esophageal, and laryngeal blistering, corneal erosions, anemia, growth retardation, and increased risk of aggressive squamous cell carcinoma in severe DEB. Diagnosis is based on clinical features, skin biopsy for immunofluorescence mapping, transmission electron microscopy, and genetic testing. Management is multidisciplinary: gentle skin care, wound dressings (non-adherent), prevention and treatment of infection, nutritional support, esophageal dilation for strictures, physical therapy for contractures, and surveillance for squamous cell carcinoma. Emerging therapies include gene therapy (beremagene geperpavec for DEB), protein replacement therapy, and cell-based therapies. Prognosis varies widely: EBS has normal life expectancy; severe JEB and recessive DEB cause significant morbidity and mortality from infection, malnutrition, and aggressive squamous cell carcinoma.

%-------------------------------------------------------------
\section{Connective Tissue Skin Diseases}
%-------------------------------------------------------------

%-------------------------------------------------------------

%-------------------------------------------------------------

%-------------------------------------------------------------

%-------------------------------------------------------------

%-------------------------------------------------------------

\label{sec:Connective Tissue Skin Diseases}
%-------------------------------------------------------------%-------------------------------------------------------------

\subsection{Dermatomyositis}
\label{sec:skin:Dermatomyositis}
Dermatomyositis is an idiopathic inflammatory myopathy with characteristic skin manifestations, affecting both children and adults. The condition results from autoimmune-mediated inflammation of skin and muscle. Cutaneous features include heliotrope rash (violaceous discoloration of the eyelids), Gottron's papules (erythematous to violaceous papules overlying the knuckles), V-sign (erythematous rash on the anterior neck and chest), shawl sign (erythematous rash over shoulders and upper back), mechanic's hands (rough, cracked skin on the fingers), and nail fold capillary changes. Proximal muscle weakness is the hallmark of myopathy. Clinical features include both skin and muscle manifestations; dermatomyositis in adults carries an increased risk of underlying malignancy (ovarian, lung, GI), necessitating age-appropriate cancer screening. Diagnosis involves CK levels, myositis-specific antibodies (anti-Mi-2, anti-MDA5, anti-TIF1-$\gamma$), muscle MRI, EMG, and muscle biopsy. Management includes corticosteroids, steroid-sparing agents (methotrexate, azathioprine, mycophenolate), IVIG, and rituximab. Prognosis is variable; malignancy screening is essential in adults.

\subsection{Systemic Lupus Erythematosus (Skin Manifestations)}
\label{sec:Systemic Lupus Erythematosus (Skin Manifestations)}
Cutaneous lupus erythematosus (CLE) encompasses a spectrum of skin involvement in systemic lupus erythematosus (SLE) or as a primary skin-limited condition. Acute cutaneous lupus presents as a malar (butterfly) rash across the cheeks and bridge of the nose, sparing the nasolabial folds, and is associated with active systemic disease. Subacute cutaneous lupus presents as annular polycyclic or papulosquamous plaques on sun-exposed areas, strongly associated with anti-Ro/SSA antibodies. Chronic cutaneous lupus (discoid lupus) presents as well-demarcated, scaly, erythematous plaques with follicular plugging and atrophic scarring, predominantly on the scalp, face, and ears. The condition results from autoimmune-mediated skin inflammation triggered by UV exposure. Diagnosis is based on clinical appearance, skin biopsy, and serological testing (ANA, anti-Ro/SSA, anti-dsDNA). Management includes sun protection, topical corticosteroids, antimalarials (hydroxychloroquine, first-line), and for refractory disease, immunosuppressants (methotrexate, mycophenolate, belimumab). Prognosis depends on the subtype and systemic involvement.

%-------------------------------------------------------------
\section{Cutaneous Infections}
%-------------------------------------------------------------

%-------------------------------------------------------------

%-------------------------------------------------------------

%-------------------------------------------------------------

%-------------------------------------------------------------

%-------------------------------------------------------------

\label{sec:Cutaneous Infections}
%-------------------------------------------------------------%-------------------------------------------------------------

\subsection{Head Lice (Pediculosis Capitis)}
\label{sec:Head Lice}
Head lice is infestation of the scalp by \textit{Pediculus humanus capitis}, common in school-age children (6--12 million cases/year in the US). Transmission occurs through direct head-to-head contact; fomite transmission is less common. The condition results from louse infestation with eggs (nits) laid at the hair base that hatch in 7--10 days and mature in 2--3 weeks. Clinical features include scalp pruritus (from hypersensitivity to louse saliva), visible nits (0.8 mm, yellow-white, firmly attached to the hair shaft), and less commonly live lice. Diagnosis is by visual inspection with a fine-toothed comb; Wood's lamp shows nits as pale blue. First-line management is permethrin 1\% or pyrethrins (OTC), applied for 10 minutes and repeated in 7--10 days; second-line options include malathion 0.5\% lotion, benzyl alcohol 5\%, spinosad 0.9\%, and ivermectin 0.5\% lotion; oral ivermectin is used for refractory cases. Wet combing is a non-chemical option. All close contacts should be checked; environmental measures include washing bedding in hot water. Prognosis is excellent with appropriate treatment; no-nit policies in schools are not recommended by the AAP/CDC.

\subsection{Impetigo}
\label{sec:Impetigo}
Impetigo is a highly contagious superficial bacterial skin infection, most common in children aged 2--5 years. The condition results from inoculation of bacteria (\textit{Staphylococcus aureus} or \textit{Streptococcus pyogenes}) into minor skin abrasions, insect bites, or areas of dermatitis. It occurs in two forms: non-bullous impetigo (70\% of cases), caused by \textit{S.~aureus} or \textit{S.~pyogenes}, presenting with honey-colored crusted erosions typically on the face and extremities; and bullous impetigo, caused by exfoliative toxin-producing \textit{S.~aureus}, presenting with flaccid, fluid-filled bullae that rupture leaving thin, brown crusts. Complications are rare but include cellulitis, lymphangitis, and post-streptococcal glomerulonephritis (from \textit{S.~pyogenes}). Diagnosis is clinical. Management for localized disease is topical mupirocin or retapamulin ointment; oral antibiotics (cephalexin, flucloxacillin, or dicloxacillin) are indicated for widespread disease, poor response to topical therapy, or lymphadenopathy. Prognosis is excellent with treatment; adequate hygiene is essential to prevent spread.

\subsection{Scabies}
\label{sec:Scabies}
Scabies is a parasitic infestation of the skin caused by the mite \textit{Sarcoptes scabiei var.~hominis}. Transmission occurs through prolonged skin-to-skin contact; fomite transmission is less common but possible. The condition results from the female mite burrowing into the stratum corneum, laying eggs and triggering a delayed type IV hypersensitivity reaction. Clinical features include intensely pruritic (especially nocturnal), papular, and vesicular eruptions in characteristic locations: interdigital web spaces, wrists, elbows, axillae, periumbilical area, waistband, buttocks, and genitalia in males; burrows (thin, linear, grayish-white tracks) are pathognomonic. Crusted (Norwegian) scabies is a highly contagious variant in immunocompromised patients with widespread, thick, hyperkeratotic crusts. Diagnosis is based on clinical examination; dermoscopy may reveal the "delta wing jet" appearance; skin scraping confirms. Management is with topical permethrin 5\% cream (applied to the entire body from neck down, left on for 8--14 hours, repeated in 1 week) or oral ivermectin (200 $\mu$g/kg, repeated in 1--2 weeks); close contacts should be treated simultaneously. Prognosis is excellent with treatment; post-scabetic pruritus may persist for weeks.

\subsection{Tinea Infections}
\label{sec:Tinea Infections}
Tinea infections (dermatophytosis) are superficial fungal infections caused by dermatophytes (Trichophyton, Microsporum, and Epidermophyton species) that infect keratinized tissue. Named by anatomic location: tinea pedis (athlete's foot, the most common, affecting the feet, particularly between toes), tinea corporis (ringworm of the body, annular, erythematous, scaly plaques with central clearing), tinea cruris (jock itch, affecting the groin and inner thighs), tinea capitis (scalp infection, most common in children, may cause alopecia and kerion formation), tinea unguium (onychomycosis, nail thickening, discoloration, and crumbling), and tinea versicolor (caused by \textit{Malassezia} species, actually a pityriasis, presenting with hypo- or hyperpigmented macules on the trunk). Risk factors include warm, moist environments, tight clothing, shared bathing facilities, immunosuppression, and diabetes. Diagnosis is based on KOH preparation showing hyphae and spores, fungal culture, dermoscopy, or Wood's lamp examination. Management depends on location and severity: topical antifungals (terbinafine, clotrimazole, miconazole) for localized tinea corporis, cruris, and pedis; oral antifungals (terbinafine, itraconazole, fluconazole) for tinea capitis, extensive tinea corporis, tinea unguium, and refractory cases. Tinea capitis requires systemic therapy due to follicular involvement. Prognosis is excellent with appropriate antifungal therapy.

\subsection{Wound Infections and Surgical Site Infections}
\label{sec:Wound Infections and Surgical Site Infections}
Wound infections and surgical site infections (SSIs) are infections occurring within 30 days of surgery (or 1 year for implant procedures) at the incision site or in deep tissues, accounting for approximately 20\% of healthcare-associated infections. Risk factors include patient factors (diabetes, obesity, immunosuppression, smoking, malnutrition) and procedural factors (contamination level, operative duration, implant placement). The condition results from bacterial inoculation during or after surgery; most SSIs are caused by \textit{S.~aureus} (including MRSA) and coagulase-negative staphylococci. SSIs are classified as superficial incisional (involving skin and subcutaneous tissue), deep incisional (involving fascia and muscle), or organ/space (involving any part opened or manipulated during surgery). Clinical features include erythema, warmth, swelling, pain, and purulent drainage from the wound. Diagnosis is clinical; wound cultures guide antibiotic therapy. Management involves wound opening and drainage for superficial SSIs, wound packing, and targeted antibiotic therapy; deep SSIs and organ/space infections may require surgical debridement or reoperation. Prognosis is generally favorable with appropriate management; prevention strategies include appropriate antibiotic prophylaxis, skin preparation with chlorhexidine-alcohol, and aseptic technique.

%-------------------------------------------------------------
\section{Eczematous Dermatoses}
%-------------------------------------------------------------

%-------------------------------------------------------------

%-------------------------------------------------------------

%-------------------------------------------------------------

%-------------------------------------------------------------

%-------------------------------------------------------------

\label{sec:Eczematous Dermatoses}
%-------------------------------------------------------------%-------------------------------------------------------------

\subsection{Atopic Dermatitis}
\label{sec:Atopic Dermatitis}
Atopic dermatitis (AD, eczema) is a chronic, relapsing inflammatory skin disease characterized by intense pruritus, xerosis, and eczematous lesions with a characteristic distribution. It affects 10--20\% of children and 3--7\% of adults. The condition results from epidermal barrier dysfunction (filaggrin mutations in 30--50\% of cases), immune dysregulation (Th2-predominant inflammation), and microbial dysbiosis (colonization with \textit{Staphylococcus aureus}); the "atopic march" links AD to food allergy, allergic rhinitis, and asthma. Clinical features vary by age: infantile AD involves the face and extensor surfaces; childhood AD affects flexural areas (antecubital and popliteal fossae); adult AD presents with lichenified plaques on flexural surfaces. Diagnosis is clinical. Management includes regular emollient use, topical corticosteroids (first-line for flares), topical calcineurin inhibitors (tacrolimus, pimecrolimus for sensitive areas), and for moderate-to-severe disease, biologic therapies (dupilumab [anti-IL-4R$\alpha$], tralokinumab [anti-IL-13]) and JAK inhibitors (abrocitinib, baricitinib). Prognosis is variable; the condition tends to improve with age.

\subsection{Contact Dermatitis}
\label{sec:Contact Dermatitis}
Contact dermatitis is inflammation of the skin caused by direct contact with an irritant or allergen. Irritant contact dermatitis (ICD) accounts for 80\% of cases and is caused by direct chemical damage to the skin (soaps, detergents, solvents, acids). Allergic contact dermatitis (ACD) is a type IV delayed hypersensitivity reaction requiring prior sensitization to the allergen. Common allergens include nickel, poison ivy/oak/urushiol, fragrances, preservatives, rubber chemicals, and hair dyes. Clinical features include pruritic, erythematous, vesicular or bullous dermatitis in the distribution of contact. Diagnosis is based on history and patch testing for ACD. Management includes identification and avoidance of the causative agent, topical corticosteroids, and in severe cases, systemic corticosteroids. Prognosis is excellent with avoidance of the offending agent.

\subsection{Nummular Dermatitis}
\label{sec:Nummular Dermatitis}
Nummular (discoid) dermatitis is a chronic inflammatory skin condition presenting as coin-shaped, well-demarcated, erythematous plaques with serous crusting, often on the extremities. It is more common in winter and in patients with dry skin. The condition results from skin barrier disruption and environmental factors. Clinical features include intense pruritus and characteristic coin-shaped lesions. Diagnosis is clinical; KOH preparation distinguishes it from tinea (negative in nummular dermatitis). Management includes potent topical corticosteroids, moisturizers, and avoidance of skin irritants. Prognosis is chronic with flares.

\subsection{Seborrheic Dermatitis}
\label{sec:Seborrheic Dermatitis}
Seborrheic dermatitis is a chronic inflammatory skin condition affecting areas rich in sebaceous glands (scalp, face, chest, back). It is associated with \textit{Malassezia} species (lipophilic yeast) and may be triggered by stress, cold weather, and neurological conditions (Parkinson's disease). In infants, it manifests as "cradle cap" (thick, yellowish scales on the scalp). Clinical features in adults include greasy, yellowish scales on an erythematous base, affecting the scalp (dandruff in mild form), eyebrows, nasolabial folds, and ears. Diagnosis is clinical. Management includes topical antifungal agents (ketoconazole shampoo/cream), low-potency topical corticosteroids, and calcineurin inhibitors. Prognosis is chronic with periodic flares.

%-------------------------------------------------------------
\section{Pigmentary Disorders}
%-------------------------------------------------------------

%-------------------------------------------------------------

%-------------------------------------------------------------

%-------------------------------------------------------------

%-------------------------------------------------------------

%-------------------------------------------------------------

\label{sec:Pigmentary Disorders}
%-------------------------------------------------------------%-------------------------------------------------------------

\subsection{Melasma}
\label{sec:Melasma}
Melasma is a common acquired hyperpigmentation disorder characterized by symmetric, irregular, brown-to-gray-brown macules and patches on sun-exposed areas, predominantly the cheeks, forehead, upper lip, and chin. It predominantly affects women with darker skin phototypes (Fitzpatrick III--VI). The condition results from UV-induced stimulation of melanocytes with increased melanin production and dermal deposition, exacerbated by sun exposure, pregnancy ("chloasma"), oral contraceptives, and hormonal therapies. Diagnosis is clinical; Wood's lamp examination helps classify depth (epidermal, dermal, mixed). Management includes strict sun protection (broad-spectrum sunscreen with iron oxide), topical depigmenting agents (hydroquinone 4\%, tretinoin, azelaic acid, tranexamic acid), and chemical peels. Prognosis is variable; recurrence is common, and combination therapy yields the best results.

\subsection{Oculocutaneous Albinism}
\label{sec:Oculocutaneous Albinism}
Oculocutaneous albinism (OCA) is a group of autosomal recessive disorders characterized by reduced or absent melanin production in the skin, hair, and eyes, with a global prevalence of 1 in 17,000. The condition results from mutations in genes involved in melanin biosynthesis: OCA1 (TYR, tyrosinase deficiency, most severe, absence of pigment), OCA2 (OCA2/P gene, common in Africans, some residual pigment), OCA3 (TYRP1, rufous albinism in Africans), and OCA4 (SLC45A2, common in Japanese). Clinical features include hypopigmentation of skin and hair ranging from white (OCA1) to light yellow or brown (OCA2/3/4), translucent irides, nystagmus, photophobia, reduced visual acuity, foveal hypoplasia, optic nerve misrouting, and strabismus. Patients have markedly increased risk of skin cancer (squamous cell carcinoma, basal cell carcinoma, melanoma) due to lack of UV protection. Diagnosis is clinical with genetic confirmation. Management involves lifelong photoprotection (broad-spectrum sunscreen, protective clothing, UV-protective eyewear), regular dermatologic surveillance for skin cancer, ophthalmologic care for visual impairment (low vision aids, tinted lenses, strabismus surgery), and skin cancer treatment. Prognosis: normal life expectancy with adequate sun protection; cutaneous malignancies are the major cause of morbidity and mortality.

\subsection{Vitiligo}
\label{sec:Vitiligo}
Vitiligo is an acquired autoimmune disorder characterized by selective destruction of melanocytes, resulting in well-demarcated, milky-white macules and patches, affecting 0.5--2\% of the global population. The condition results from autoimmune-mediated melanocyte destruction. It may be segmental (unilateral, dermatomal, often with early onset and stable course) or non-segmental (generalized, bilateral, with variable progression). Loss of pigment may occur in any area but commonly affects the face, hands, elbows, knees, and genitalia. Vitiligo is associated with other autoimmune diseases (thyroid disease, type 1 diabetes, Addison's disease, pernicious anemia); the Koebner phenomenon (new lesions at sites of trauma) may occur. Diagnosis is clinical, supported by Wood's lamp examination. Management includes topical corticosteroids, topical calcineurin inhibitors (tacrolimus, especially for facial lesions), narrowband UVB phototherapy, and for refractory localized disease, surgical melanocyte transplantation; ruxolitinib cream (JAK1/2 inhibitor) has been approved for non-segmental vitiligo. Sun protection is essential. Prognosis is variable; repigmentation is achievable but recurrence may occur.

%-------------------------------------------------------------
\section{Psoriasis}
%-------------------------------------------------------------

%-------------------------------------------------------------

%-------------------------------------------------------------

%-------------------------------------------------------------

%-------------------------------------------------------------

%-------------------------------------------------------------

\label{sec:section:Psoriasis}
Psoriasis is a chronic, immune-mediated inflammatory dermatosis affecting skin, nails, and joints in susceptible individuals. Dysregulated Th17 and Th1 cytokine signaling stimulates rapid keratinocyte hyperproliferation and epidermal hyperplasia. Clinical management follows a stepwise approach utilizing topical agents, phototherapy, conventional oral immunosuppressants, and targeted biologic therapies.

%-------------------------------------------------------------%-------------------------------------------------------------

\subsection{Psoriasis}
\label{sec:Psoriasis}

Psoriasis is a chronic, immune-mediated inflammatory skin disease characterized by well-demarcated, erythematous plaques with silvery-white scales, affecting 2--3\% of the global population. The condition results from dysregulated Th1/Th17 immune responses leading to keratinocyte hyperproliferation (epidermal turnover time reduced from 28 days to 3--5 days). Plaque psoriasis (psoriasis vulgaris) is the most common form, presenting with symmetrically distributed plaques on extensor surfaces (elbows, knees), scalp, and lower back. Other forms include guttate (small, drop-shaped lesions, often post-streptococcal), inverse (flexural areas), pustular, and erythrodermic psoriasis. Nail changes (pitting, onycholysis, oil-drop discoloration) are common; up to 30\% of patients develop psoriatic arthritis. Diagnosis is clinical. Management follows a stepwise approach based on severity: topical therapies (corticosteroids, vitamin D analogues [calcipotriene], topical calcineurin inhibitors) for mild disease; phototherapy (narrowband UVB) for moderate disease; and systemic therapy for moderate-to-severe disease: conventional agents (methotrexate, cyclosporine, acitretin), and biologics targeting specific cytokines (TNF inhibitors, IL-17 inhibitors, IL-23 inhibitors, IL-12/23 inhibitor [ustekinumab]). Prognosis is chronic but can be well-controlled with appropriate therapy.

\subsection{Psoriatic Arthritis}
\label{sec:skin:Psoriatic Arthritis}
Psoriatic arthritis (PsA) is an inflammatory arthropathy associated with psoriasis, affecting 30\% of patients with cutaneous psoriasis. It is a seronegative spondyloarthropathy (negative for rheumatoid factor and anti-CCP antibodies). The condition results from immune-mediated joint inflammation. Five clinical patterns exist: oligoarthritis (most common), polyarthritis, arthritis mutilans (severe, destructive), spondylitis, and dactylitis ("sausage digits"). Extra-articular manifestations include enthesitis (Achilles, plantar fascia), nail changes, eye inflammation, and cardiovascular disease. Diagnosis is based on the CASPAR criteria. Management includes NSAIDs, DMARDs (methotrexate, sulfasalazine), and biologics (TNF inhibitors, IL-17 inhibitors, IL-23 inhibitors, JAK inhibitors, PDE4 inhibitors). Prognosis is improved with early diagnosis and treatment to prevent irreversible joint damage.

%-------------------------------------------------------------
\section{Skin Cancer}
%-------------------------------------------------------------

%-------------------------------------------------------------

Skin cancer represents the most common malignancy in humans, categorized into non-melanoma skin cancers (basal cell carcinoma and squamous cell carcinoma) and cutaneous melanoma. Ultraviolet radiation exposure induces DNA damage and oncogenic mutations in keratinocytes or melanocytes. Early detection, precise dermatopathological staging, surgical excision, and modern immunotherapy achieve excellent long-term survival.

%-------------------------------------------------------------

%-------------------------------------------------------------

%-------------------------------------------------------------

%-------------------------------------------------------------

\label{sec:Skin Cancer}
%-------------------------------------------------------------%-------------------------------------------------------------

\subsection{Actinic Keratosis}
\label{sec:Actinic Keratosis}
Actinic keratosis is a precancerous lesion caused by chronic UV radiation exposure, representing the earliest stage of the squamous cell carcinoma spectrum. The condition results from UV-induced keratinocyte dysplasia. Clinical features include rough, scaly, erythematous papules or plaques on sun-exposed skin (face, scalp, ears, forearms, hands); AKs carry a risk of progression to invasive squamous cell carcinoma (approximately 1\% per lesion per year). Diagnosis is clinical; biopsy may be required to exclude carcinoma. Management options depend on the number and location of lesions: cryotherapy (liquid nitrogen), topical 5-fluorouracil, imiquimod, ingenol mebutate, photodynamic therapy, and field therapy for extensive actinic damage. Prognosis is excellent with treatment; regular skin surveillance is recommended.

\subsection{Basal Cell Carcinoma}
\label{sec:Basal Cell Carcinoma}
Basal cell carcinoma (BCC) is the most common skin cancer worldwide, arising from basal cells of the epidermis. It is strongly associated with cumulative ultraviolet radiation exposure and predominantly affects fair-skinned individuals. BCC is locally invasive but rarely metastasizes. Clinical subtypes include nodular (most common, pearly papule with telangiectasia and central ulceration---"rodent ulcer"), superficial (erythematous, scaly plaque, often on the trunk), and morpheaform (sclerosing, whitish, scar-like). Diagnosis is based on biopsy. Management depends on subtype, size, and location: surgical excision, Mohs micrographic surgery (for high-risk locations such as the face, with highest cure rate of 99\%), electrodesiccation and curettage (for superficial, low-risk lesions), and topical therapy (5-fluorouracil, imiquimod) for superficial BCC; vismodegib and sonidegib (Hedgehog pathway inhibitors) are used for advanced or metastatic BCC. Prognosis is excellent; metastasis is extremely rare.

\subsection{Melanoma}
\label{sec:Melanoma}
Melanoma is the most lethal skin cancer, arising from melanocytes. Although representing only 1\% of skin cancers, it accounts for the majority of skin cancer deaths. Risk factors include UV exposure (intermittent intense exposure and sunburns), fair skin, multiple nevi, atypical nevi, family history, and genetic susceptibility (CDKN2A, CDK4 mutations). The ABCDE criteria aid in identification: Asymmetry, Border irregularity, Color variation, Diameter >6 mm, and Evolution. The four main subtypes are superficial spreading (most common), nodular (most aggressive, may not follow ABCDE), lentigo maligna (on chronically sun-damaged skin), and acral lentiginous (palms, soles, nail beds; most common melanoma in people of color). The condition results from malignant transformation of melanocytes driven by UV-induced and genetic mutations (BRAF V600, NRAS, KIT). Diagnosis requires excisional biopsy with adequate margins; staging uses Breslow thickness, ulceration, mitotic rate, and lymph node status. Management of thin melanoma ($\leq$1 mm) is surgical excision with sentinel lymph node biopsy for those $\geq$0.8 mm or with ulceration; immunotherapy (anti-PD-1 [nivolumab, pembrolizumab], anti-CTLA-4 [ipilimumab], combination therapy) and targeted therapy (BRAF/MEK inhibitors for BRAF V600-mutated melanoma) have revolutionized treatment of advanced melanoma. Prognosis depends on stage at diagnosis; early detection is critical.

\subsection{Squamous Cell Carcinoma}
\label{sec:Squamous Cell Carcinoma}
Cutaneous squamous cell carcinoma (cSCC) is the second most common skin cancer, arising from keratinocytes in the epidermis. Risk factors include cumulative UV exposure, fair skin, immunosuppression (organ transplant recipients have 65-fold increased risk), chronic wounds, and actinic keratoses (precursor lesions). The condition results from UV-induced DNA damage and immunosuppression leading to malignant transformation of keratinocytes. cSCC is more aggressive than BCC, with metastatic potential (2--5\% for primary lesions, higher for recurrent or immunosuppressed patients). Clinical features vary from an erythematous, scaly plaque to a firm, ulcerated nodule. Diagnosis is based on biopsy. Management includes surgical excision, Mohs surgery for high-risk lesions, and radiation therapy for inoperable cases; cemiplimab (anti-PD-1) is approved for advanced cSCC. Prognosis is generally favorable but worse than BCC; metastatic risk is higher in immunocompromised patients.

%-------------------------------------------------------------
\section{Skin Infections}
%-------------------------------------------------------------

%-------------------------------------------------------------

Skin infections comprise a wide range of bacterial, viral, fungal, and parasitic cutaneous invasions affecting the epidermis, dermis, or subcutaneous tissue. Disruption of the epidermal barrier allows opportunistic pathogen entry, triggering localized erythema, edema, purulence, or systemic sepsis. Prompt targeted antimicrobial therapy and surgical drainage of localized abscesses remain fundamental to clinical resolution.

%-------------------------------------------------------------

%-------------------------------------------------------------

%-------------------------------------------------------------

%-------------------------------------------------------------

\label{sec:Skin Infections}
%-------------------------------------------------------------%-------------------------------------------------------------

\subsection{Abscess}
\label{sec:Abscess}
A cutaneous abscess is a localized collection of pus in the dermis or subcutaneous tissue, often caused by \textit{S.~aureus} (including MRSA). The condition results from infection of a hair follicle (folliculitis progressing to furuncle to carbuncle), blocked sweat gland, or direct inoculation. Clinical features include a painful, fluctuant, erythematous nodule with surrounding induration. Diagnosis is clinical. Management is incision and drainage (I\&D), which is the primary intervention; antibiotics are not routinely required for simple abscesses after adequate I\&D but are indicated for surrounding cellulitis, immunocompromised patients, multiple lesions, or facial location. Prognosis is excellent with adequate drainage; recurrent abscesses should prompt evaluation for risk factors (diabetes, HIV, nasal \textit{S.~aureus} carriage, community MRSA).

\subsection{Cellulitis}
\label{sec:Cellulitis}
Cellulitis is a bacterial infection of the dermis and subcutaneous tissue, most commonly caused by \textit{Streptococcus pyogenes} (group A strep) and \textit{Staphylococcus aureus}. Predisposing factors include skin barrier disruption (fissures, insect bites, surgical wounds), edema, chronic venous insufficiency, and immunosuppression. The condition results from bacterial entry through disrupted skin barriers. Clinical features include a well-demarcated area of erythema, warmth, swelling, tenderness, and sometimes lymphangitis and regional lymphadenopathy; the lower extremities are most commonly affected. Complications include abscess formation, necrotizing fasciitis (a surgical emergency), bacteremia, and sepsis. Diagnosis is clinical; wound cultures and blood cultures may guide therapy in severe cases. Management includes oral antibiotics (dicloxacillin, cephalexin, or clindamycin for MRSA coverage) for mild cases and IV antibiotics (ceftriaxone, vancomycin, or daptomycin) for severe disease. Prognosis is excellent with appropriate antibiotic therapy; recurrent cellulitis may require prophylactic antibiotics.

\subsection{Herpes Simplex}
\label{sec:Herpes Simplex}
Herpes simplex virus (HSV) infection causes recurrent, painful vesicular eruptions on an erythematous base. HSV-1 typically causes orolabial herpes (cold sores), while HSV-2 typically causes genital herpes, though there is increasing overlap. The condition results from HSV infection with establishment of latency in sensory ganglia (trigeminal for HSV-1, sacral for HSV-2) and reactivation with stress, illness, UV exposure, or immunosuppression. Primary infection may be asymptomatic or present with fever, malaise, and widespread vesicles (especially in children---gingivostomatitis). Recurrent disease is milder, with clustered vesicles at a consistent site, lasting 7--10 days. Diagnosis is based on clinical examination, viral culture, PCR, or direct fluorescent antibody testing. Management includes topical or oral antivirals (acyclovir, valacyclovir, famciclovir), most effective when started early; suppressive antiviral therapy reduces recurrence frequency. Prognosis is good; recurrences decrease over time in many patients.

\subsection{Herpes Zoster}
\label{sec:Herpes Zoster}
Herpes zoster (shingles) results from reactivation of latent varicella-zoster virus (VZV) in dorsal root or cranial nerve ganglia, typically occurring decades after primary chickenpox infection. The condition results from declining cell-mediated immunity allowing VZV reactivation. Clinical features include a unilateral, dermatomal distribution of painful vesicular eruption, most commonly affecting the thoracic dermatomes; prodromal pain may precede the rash by 1--3 days. Complications include postherpetic neuralgia (PHN, persistent pain for >3 months after rash, most common in patients >50), bacterial superinfection, motor neuropathy, and disseminated zoster in immunocompromised patients; ophthalmic zoster (V1 distribution) requires urgent ophthalmologic evaluation. Diagnosis is clinical; PCR or DFA testing may confirm atypical cases. Management with antivirals (valacyclovir or famciclovir) should be initiated within 72 hours of rash onset. Prognosis is generally favorable; the Shingrix recombinant zoster vaccine is recommended for adults $\geq$50 to prevent herpes zoster and PHN.

\subsection{Molluscum Contagiosum}
\label{sec:Molluscum Contagiosum}
Molluscum contagiosum is a viral skin infection caused by a poxvirus, transmitted by direct contact (including sexual contact) or fomites. The condition results from poxvirus infection of epidermal keratinocytes. Clinical features include 2--5 mm, dome-shaped, flesh-colored papules with central umbilication ("dell"), often with a white waxy core; lesions may be solitary or numerous and can be spread by scratching (autoinoculation). In immunocompromised patients (HIV/AIDS), molluscum can be widespread, large, and refractory to treatment. Diagnosis is clinical; dermoscopy shows a central pore with white amorphous material. Management options include cryotherapy, curettage, topical agents (cantharidin, imiquimod, potassium hydroxide), and mechanical extraction. Prognosis is excellent; self-resolution occurs in immunocompetent individuals within months to years.

\subsection{Warts}
\label{sec:Warts}
Warts (verrucae) are benign epithelial proliferations caused by human papillomavirus (HPV) infection, with over 100 HPV types identified. The condition results from HPV infection of keratinocytes causing hyperkeratosis. Common warts (verruca vulgaris) present as rough, hyperkeratotic papules, typically on the hands and fingers. Plantar warts (verruca plantaris) occur on the soles of the feet, growing inward due to pressure, causing pain with walking. Flat warts (verruca plana) are smooth, flat-topped papules on the face and dorsum of hands. Genital warts (condylomata acuminata) are caused by HPV types 6 and 11. Diagnosis is clinical. Management options include salicylic acid (first-line for common and plantar warts), cryotherapy, cantharidin, imiquimod (for genital warts), and surgical destruction for refractory lesions. Prognosis is excellent; many warts resolve spontaneously within 1--2 years due to cell-mediated immunity.

%-------------------------------------------------------------
\section{Urticaria and Related Disorders}
%-------------------------------------------------------------

%-------------------------------------------------------------

%-------------------------------------------------------------

%-------------------------------------------------------------

%-------------------------------------------------------------

%-------------------------------------------------------------

\label{sec:Urticaria and Related Disorders}
%-------------------------------------------------------------

\subsection{Alopecia Areata}
\label{sec:Alopecia Areata}
Alopecia areata is an autoimmune dermatological disorder characterized by non-scarring patchy hair loss affecting hair-bearing areas of the body, most commonly the scalp and beard area. The condition results from CD8+ T cell-mediated autoimmune attack directed against hair follicles in the anagen (growth) phase, with collapse of hair follicle immune privilege. Clinical features include sudden onset of smooth, round or oval, circumscribed patches of hair loss without erythema or scaling, featuring pathognomonic 'exclamation mark' hairs (short broken hairs tapering near the scalp surface) at active patch margins. Severe variants include alopecia totalis (complete loss of scalp hair) and alopecia universalis (complete loss of all scalp and body hair). Nail changes (pitting, trachyonychia) occur in 10--20\% of patients. Diagnosis is clinical based on characteristic morphology and dermoscopy (trichoscopy). Management includes topical or intralesional corticosteroid injections (triamcinolone acetonide) for localized disease, topical minoxidil, immunotherapy, and oral Janus kinase (JAK) inhibitors (baricitinib, ritlecitinib) for severe diffuse disease. Prognosis is variable; spontaneous regrowth occurs in localized disease, but relapse is common.

\subsection{Pemphigus Vulgaris}
\label{sec:Pemphigus Vulgaris}
Pemphigus vulgaris (PV) is a potentially life-threatening autoimmune blistering disease characterized by IgG autoantibodies against desmoglein 3 (and sometimes desmoglein 1), components of desmosomes that maintain keratinocyte adhesion. The condition results from autoantibody-mediated loss of keratinocyte adhesion leading to intraepidermal blisters (suprabasal acantholysis). Clinical features include mucosal involvement (oral erosions) in 50--70\% of cases, often preceding cutaneous lesions; skin blisters are flaccid, easily ruptured (positive Nikolsky sign), and heal without scarring. Diagnosis requires skin biopsy showing suprabasal acantholysis, direct immunofluorescence showing intercellular IgG and C3 deposition ("chicken-wire" pattern), and detection of anti-desmoglein antibodies by ELISA. Management requires immunosuppression: systemic corticosteroids combined with steroid-sparing agents (azathioprine, mycophenolate, rituximab); rituximab (anti-CD20) has become a first-line treatment, often achieving drug-free remission. Prognosis has improved significantly with modern immunosuppressive therapy.

\subsection{Pilonidal Sinus}
\label{sec:Pilonidal Sinus}
Pilonidal sinus (pilonidal disease) is a chronic inflammatory condition affecting the skin and subcutaneous tissue of the intergluteal cleft (natal cleft), characterized by hair-containing sinus tracts. The condition results from mechanical occlusion of hair follicles in the sacrococcygeal region, leading to follicle rupture, foreign-body granulomatous reaction to entrapped shed hairs, and secondary bacterial infection (Staphylococcus aureus, anaerobes). Risk factors include male sex, hirsutism, deep natal cleft, sedentary lifestyle, and obesity. Clinical features range from an asymptomatic small pit in the natal cleft to an acute pilonidal abscess (severe localized pain, erythema, fluctuating mass, fever) or chronic purulent sinus drainage. Diagnosis is clinical. Management of acute abscess requires prompt incision and drainage; chronic sinus disease is treated with surgical excision or unroofing (lay-open) with primary closure or secondary healing, accompanied by local hair removal.

\subsection{Urticaria}
\label{sec:Urticaria}
Urticaria is a common skin condition characterized by transient, itchy, erythematous wheals (hives) with central clearing, affecting 15--25\% of the population at some point in life. Acute urticaria (duration <6 weeks) is usually caused by IgE-mediated reactions to foods (shellfish, nuts, eggs), medications (antibiotics, NSAIDs), infections (viral URIs), insect stings, or physical stimuli. Chronic urticaria (duration >6 weeks) is idiopathic in 50\% and autoimmune in 30--50\% (IgG autoantibodies against Fc$\varepsilon$RI or IgE). The condition results from mast cell degranulation releasing histamine and other vasoactive mediators, causing dermal edema and vasodilation. Clinical features include well-circumscribed, raised, erythematous wheals with intense pruritus, ranging from a few millimeters to several centimeters. Individual wheals resolve within 24 hours without scarring. Angioedema accompanies 40\% of cases. Diagnosis is by history and physical examination; chronic urticaria requires CBC, ESR, TSH, and autoantibody testing. Management includes first-generation and second-generation H1 antihistamines (cetirizine, fexofenadine, loratadine) as first-line therapy; H2 antihistamines (ranitidine, famotidine) may be added. For refractory chronic urticaria: omalizumab (anti-IgE monoclonal antibody, highly effective), leukotriene receptor antagonists, cyclosporine, or hydroxychloroquine. Short courses of oral corticosteroids for severe exacerbations. Prognosis: acute urticaria resolves with trigger avoidance; chronic urticaria resolves within 1 year in 50\% and within 5 years in 80\%.

